\appendix

\section{Deferred Proofs}
\label{app:proofs}

This appendix collects detailed proofs of results stated in
the main text.  Each proof is presented with intermediate steps
to make the arguments self-contained.

\subsection*{Proof of Theorem~\ref{thm:chebyshev} (Chebyshev States)}

\begin{proof}
	For real signals ($\theta = 0$), the stereographic signal
	operator reduces to the rotation~\eqref{eq:rotation-form}:
	\begin{equation}
		S_z\big|_{\theta=0} = R(\phi) = \begin{pmatrix} \cos\phi & -\sin\phi \\ \sin\phi & \cos\phi \end{pmatrix},
		\tag{A1}
	\end{equation}
	with $\phi = \arccos(\rtilde)$ and
	$\rtilde = r/\sqrt{1+r^2}$.  Since $R(\phi)$ is a
	rotation matrix in $\mathrm{SO}(2)$, its $k$-th power is simply
	$R(k\phi)$:
	\begin{equation}
		S_z^k\ket{0} = R(k\phi)\begin{pmatrix}1\\0\end{pmatrix} = \begin{pmatrix} \cos(k\phi) \\ \sin(k\phi) \end{pmatrix}.
		\tag{A2}
	\end{equation}
	It remains to express $\cos(k\phi)$ and $\sin(k\phi)$
	in terms of $\rtilde$.  By definition of the Chebyshev
	polynomials,
	\begin{align}
		T_k(\cos\phi)             & = \cos(k\phi), \tag{A3a} \\
		U_{k-1}(\cos\phi)\sin\phi & = \sin(k\phi), \tag{A3b}
	\end{align}
	where $T_k$ is the Chebyshev polynomial of the first kind
	and $U_{k-1}$ is the Chebyshev polynomial of the second kind.
	The identity~(A3b) is the standard relation between
	Chebyshev polynomials of the second kind and trigonometric
	functions (see, e.g., Ref.~\cite{BenentiCasatiStrini2004}).

	Substituting $\cos\phi = \rtilde$ into~(A3a) gives the
	$\ket{0}$-component directly:
	$\cos(k\phi) = T_k(\rtilde)$.  For the
	$\ket{1}$-component, we use~(A3b) with
	$\sin\phi = \sqrt{1 - \rtilde^2} = \sqrt{1 - r^2/(1+r^2)} =
		1/\sqrt{1+r^2}$, so that
	\begin{equation}
		\sin(k\phi) = U_{k-1}(\rtilde)\sin\phi = \frac{U_{k-1}(\rtilde)}{\sqrt{1+r^2}}.
		\tag{A4}
	\end{equation}
	Combining (A2), (A3a), and (A4):
	\begin{equation}
		S_z^k\ket{0} = \begin{pmatrix} T_k(\rtilde) \\ U_{k-1}(\rtilde)/\sqrt{1+r^2} \end{pmatrix}
		= T_k(\rtilde)\,\ket{0} + \frac{U_{k-1}(\rtilde)}{\sqrt{1+r^2}}\,\ket{1},
		\tag{A5}
	\end{equation}
	which is the claimed result.  One can verify the
	normalization: $|T_k(\rtilde)|^2 +
		|U_{k-1}(\rtilde)|^2/(1+r^2) = \cos^2(k\phi) +
		\sin^2(k\phi) = 1$, confirming normalization.
\end{proof}

\subsection*{Proof of Corollary~\ref{cor:existence} (Existence of Phases)}

\begin{proof}
	We must show that for any polynomial pair $(P, Q)$ of
	degree $\leq d$ satisfying $|P(\rtilde)|^2 + |Q(\rtilde)|^2 = 1$
	and definite parity, there exist phases $\bm{\phi}$ such
	that $W_{\bm{\phi}}^{(\mathrm{stereo})}(r)\ket{0} =
		P(\rtilde)\ket{0} + Q(\rtilde)\ket{1}$.  The argument
	has two steps.

	\textit{Step 1: Invoke standard QSP.}
	The polynomial $P$ satisfies $|P(a)| \leq 1$ for all
	$a \in [-1,1]$ (since $|P|^2 \leq |P|^2 + |Q|^2 = 1$)
	and has definite parity by hypothesis.  By
	Theorem~\ref{thm:standard-qsp}, there exist phases
	$\bm{\phi} = (\phi_0, \ldots, \phi_d)$ such that the
	standard QSP product satisfies
	$\bra{0}U_{\bm{\phi}}^{(\mathrm{std})}(a)\ket{0} = P(a)$
	for all $a \in [-1,1]$.  In particular, this holds at
	$a = \rtilde$:
	\begin{equation}
		\bra{0}U_{\bm{\phi}}^{(\mathrm{std})}(\rtilde)\ket{0} = P(\rtilde).
		\tag{A6}
	\end{equation}

	\textit{Step 2: Transfer via the basis change.}
	By Proposition~\ref{prop:basis-change}, the stereographic
	and standard QSP products are related by the conjugation
	$W_{\bm{\phi}}^{(\mathrm{stereo})}(r)
		= V^{-1}\,U_{\bm{\phi}}^{(\mathrm{std})}(\rtilde)\,V$,
	where $V = \mathrm{diag}(1, i)$ is a fixed unitary.
	Since $V\ket{0} = \ket{0}$ and
	$\bra{0}V^{-1} = \bra{0}$ (the matrix $V$ acts trivially
	on the first basis vector), the $(0,0)$ matrix element is
	preserved:
	\begin{align}
		&\bra{0}W_{\bm{\phi}}^{(\mathrm{stereo})}(r)\ket{0}
		\nonumber\\
		&\quad= \bra{0}V^{-1}\,U_{\bm{\phi}}^{(\mathrm{std})}(\rtilde)\,V\ket{0}
		= \bra{0}U_{\bm{\phi}}^{(\mathrm{std})}(\rtilde)\ket{0}
		= P(\rtilde).
		\tag{A7}
	\end{align}
	Therefore the \emph{same phases} $\bm{\phi}$ that
	realize $P$ in the standard protocol also realize $P$
	in the stereographic protocol.  The complementary
	polynomial $Q$ is then determined by unitarity:
	$|P(\rtilde)|^2 + |Q(\rtilde)|^2 = 1$ fixes $|Q|$, and
	the phase of $Q$ is inherited from the full unitary
	matrix.
\end{proof}


\subsection*{Proof of Proposition~\ref{prop:zeros-poles} (Zeros and Poles)}

\begin{proof}
	The function $z_k(r) = \cot(k\arctan(1/r))$ has zeros
	where $\cot(\cdot) = 0$ and poles where $\cot(\cdot)$
	diverges.  The cotangent function $\cot(\alpha)$ vanishes
	when $\alpha = (2m-1)\pi/2$ (odd multiples of $\pi/2$)
	and diverges when $\alpha = m\pi$ (multiples of~$\pi$),
	for integer~$m$.  Applying this to
	$\alpha = k\arctan(1/r)$:

	\textit{Zeros.}  Setting $k\arctan(1/r) = (2m-1)\pi/2$
	gives $\arctan(1/r) = (2m-1)\pi/(2k)$, so
	$1/r = \tan\!\big((2m-1)\pi/(2k)\big)$ and
	\begin{equation}
		r_m^{(0)} = \cot\!\bigl((2m-1)\pi/(2k)\bigr),
		\quad m = 1, 2, \ldots
		\tag{A8}
	\end{equation}
	The zero $r_m^{(0)}$ lies on $r > 0$ when the argument
	$(2m-1)\pi/(2k) \in (0, \pi/2)$, i.e., when
	$m = 1, \ldots, \lfloor k/2 \rfloor$.  Larger values
	of~$m$ either give $r \leq 0$ or repeat existing zeros
	(by periodicity of cotangent).

	\textit{Poles.}  Setting $k\arctan(1/r) = m\pi$ gives
	$\arctan(1/r) = m\pi/k$, so
	\begin{equation}
		r_m^{(\infty)} = \cot\!\bigl(m\pi/k\bigr),
		\quad m = 1, 2, \ldots
		\tag{A9}
	\end{equation}
	The pole lies on $r > 0$ when $m\pi/k \in (0, \pi/2)$,
	i.e., when $m = 1, \ldots, \lfloor (k-1)/2 \rfloor$.

	\textit{Behavior at $r = 0$.}  As $r \to 0^+$,
	$\arctan(1/r) \to \pi/2$, so the argument of the
	cotangent approaches $k\pi/2$.  There are two cases:
	\begin{itemize}
		\item If $k$ is odd, $k\pi/2$ is an odd multiple
		      of $\pi/2$, so $\cot(k\pi/2) = 0$ and
		      $z_k(0^+) = 0$---a zero.
		\item If $k$ is even, $k\pi/2$ is a multiple
		      of~$\pi$, so $\cot(k\pi/2)$ diverges and
		      $z_k$ has a pole at $r = 0$.
	\end{itemize}
	This parity-dependent behavior reflects the fact
	that $T_k(0) = 0$ for odd~$k$ (numerator vanishes)
	and $U_{k-1}(0) = 0$ for even~$k$ (denominator
	vanishes), when expressed in the Chebyshev
	representation $z_k = T_k(\rtilde)\sqrt{1+r^2}/U_{k-1}(\rtilde)$.

	\textit{Total count.}  On $r > 0$, the total number of
	zeros and poles is $\lfloor k/2 \rfloor + \lfloor (k-1)/2 \rfloor$,
	which equals $k - 1$ for both odd and even~$k$.  Including
	the zero or pole at $r = 0$ brings the total to exactly~$k$
	points (counting multiplicities), consistent with the degree
	of the rational function.
\end{proof}


\subsection*{Proof of Theorem~\ref{thm:unnormalized} (Unnormalized Formulation)}

\begin{proof}
	We prove by induction on~$d$ that the product
	$\widetilde{W}_{\bm{\phi}}(r) =
		\exp[i\phi_0\PauliZ]\prod_{j=1}^{d}O_r\exp[i\phi_j\PauliZ]$
	has the form
	$\bigl(\begin{smallmatrix} \tilde{P} & -\tilde{Q}^* \\
				\tilde{Q} & \tilde{P}^* \end{smallmatrix}\bigr)$ with
	$|\tilde{P}|^2 + |\tilde{Q}|^2 = (1+r^2)^d$,
	$\deg(\tilde{P}) \leq d$, $\deg(\tilde{Q}) \leq d-1$,
	and the stated parity conditions.

	\textit{Base case ($d = 0$).}
	When $d = 0$ there are no signal operators, and the
	product is simply a phase gate:
	\begin{equation}
		\widetilde{W}_{\bm{\phi}} = \exp[i\phi_0\PauliZ] =
		\begin{pmatrix} e^{i\phi_0} & 0            \\
                0           & e^{-i\phi_0}\end{pmatrix}.
		\tag{A10}
	\end{equation}
	Reading off $\tilde{P} = e^{i\phi_0}$ and
	$\tilde{Q} = 0$, we verify:
	$|\tilde{P}|^2 + |\tilde{Q}|^2 = 1 = (1+r^2)^0$;
	$\deg(\tilde{P}) = 0 \leq 0$;
	$\deg(\tilde{Q}) = -\infty \leq -1$;
	and both satisfy trivial parity conditions (constants
	have even parity, consistent with $d = 0$ being even).

	\textit{Inductive step.}
	Suppose after $d-1$ signal applications we have
	polynomials $P, Q \in \mathbb{C}[r]$ satisfying the
	inductive hypotheses:
	\begin{enumerate}
		\item $|P|^2 + |Q|^2 = (1+r^2)^{d-1}$,
		\item $\deg(P) \leq d-1$ and $\deg(Q) \leq d-2$,
		\item $P$ has parity $(d-1)\bmod 2$ and $Q$ has
		      parity $(d-2)\bmod 2$.
	\end{enumerate}
	The product up to step $d-1$ is
	$M_{d-1} = \bigl(\begin{smallmatrix} P & -Q^* \\
				Q & P^* \end{smallmatrix}\bigr)$.
	Appending one more factor of $\exp[i\phi_d\PauliZ]\,O_r$
	from the left (recall that in the QSP product, the
	phase--signal pairs are applied in sequence) gives the
	new product matrix
	\begin{equation}
		M_d = \exp[i\phi\PauliZ]\,O_r\,M_{d-1},
		\tag{A11}
	\end{equation}
	where we write $\phi = \phi_d$ for brevity.  We compute
	this in two stages.

	\textit{Stage 1: Multiply by $O_r$.}
	The unnormalized signal operator is
	$O_r = \bigl(\begin{smallmatrix} r & -1 \\ 1 & r
			\end{smallmatrix}\bigr)$, so
	\begin{align}
		O_r\,M_{d-1}
		&= \begin{pmatrix} r & -1 \\ 1 & r \end{pmatrix}
		\begin{pmatrix} P & -Q^* \\ Q & P^* \end{pmatrix}
		\nonumber\\
		&= \begin{pmatrix} rP - Q & -(rQ^* + P^*) \\
                P + rQ & rP^* - Q^*\end{pmatrix}.
		\tag{A12}
	\end{align}
	Identifying the new polynomials as
	$\hat{P} = rP - Q$ and $\hat{Q} = P + rQ$, we verify
	the $\SU$-like structure:
	$(O_r M_{d-1})_{12} = -(rQ^* + P^*) = -\hat{Q}^*$ and
	$(O_r M_{d-1})_{22} = rP^* - Q^* = \hat{P}^*$,
	confirming the matrix has the form
	$\bigl(\begin{smallmatrix} \hat{P} & -\hat{Q}^* \\
				\hat{Q} & \hat{P}^* \end{smallmatrix}\bigr)$.

	\textit{Stage 2: Multiply by the phase gate.}
	\begin{equation}
		\exp[i\phi\PauliZ]\begin{pmatrix} \hat{P} & -\hat{Q}^* \\ \hat{Q} & \hat{P}^* \end{pmatrix}
		= \begin{pmatrix} e^{i\phi}\hat{P} & -e^{i\phi}\hat{Q}^* \\ e^{-i\phi}\hat{Q} & e^{-i\phi}\hat{P}^* \end{pmatrix}.
		\tag{A13}
	\end{equation}
	Setting $\widetilde{P} = e^{i\phi}\hat{P}
		= e^{i\phi}(rP - Q)$ and
	$\widetilde{Q} = e^{-i\phi}\hat{Q}
		= e^{-i\phi}(P + rQ)$, the new matrix retains the
	required form with $|\widetilde{P}| = |\hat{P}|$ and
	$|\widetilde{Q}| = |\hat{Q}|$ (the phase factors do not
	affect moduli).

	\textit{Verifying the norm condition.}
	We must show
	$|\widetilde{P}|^2 + |\widetilde{Q}|^2 = (1+r^2)^d$.
	Expanding:
	\begin{align}
		|\widetilde{P}|^2 + |\widetilde{Q}|^2
		 & = |rP - Q|^2 + |P + rQ|^2 \nonumber                            \\
		 & = \bigl(r^2|P|^2 - rPQ^* - rP^*Q + |Q|^2\bigr) \nonumber       \\
		 & \quad + \bigl(|P|^2 + rPQ^* + rP^*Q + r^2|Q|^2\bigr) \nonumber \\
		 & = (r^2+1)|P|^2 + (1+r^2)|Q|^2 \nonumber                        \\
		 & = (1+r^2)\bigl(|P|^2+|Q|^2\bigr) \nonumber                     \\
		 & = (1+r^2)\cdot(1+r^2)^{d-1} = (1+r^2)^d.
		\tag{A14}
	\end{align}
	The key cancellation occurs in the cross terms: the
	term $-rPQ^* - rP^*Q = -2r\,\mathrm{Re}(PQ^*)$
	from $|rP - Q|^2$ is exactly cancelled by the
	term $+rPQ^* + rP^*Q = +2r\,\mathrm{Re}(PQ^*)$
	from $|P + rQ|^2$.  This is not a coincidence---it
	reflects the orthogonality of the two columns of $O_r$:
	\begin{equation}
		\begin{pmatrix}r \\ 1\end{pmatrix}^{\!\top}
		\begin{pmatrix}-1 \\ r\end{pmatrix} = -r + r = 0.
		\tag{A15}
	\end{equation}

	\textit{Verifying the degree.}
	From $\widetilde{P} = e^{i\phi}(rP - Q)$ with
	$\deg(P) \leq d-1$ and $\deg(Q) \leq d-2$: the term
	$rP$ has degree $\leq d$ and the term $Q$ has degree
	$\leq d-2$, so $\deg(\widetilde{P}) \leq d$.
	From $\widetilde{Q} = e^{-i\phi}(P + rQ)$ with
	$\deg(P) \leq d-1$ and $\deg(rQ) \leq d-1$: both
	terms have degree $\leq d-1$, so
	$\deg(\widetilde{Q}) \leq d-1$.

	\textit{Verifying the parity.}
	Recall that a polynomial $f(r)$ has parity $p$ if
	$f(-r) = (-1)^p f(r)$.  By the inductive hypothesis,
	$P$ has parity $(d-1)\bmod 2$ and $Q$ has parity
	$(d-2)\bmod 2$.  Multiplication by $r$ flips parity
	(since $r$ is odd), so $rP$ has parity $d\bmod 2$ and
	$rQ$ has parity $(d-1)\bmod 2$.  Then:
	\begin{itemize}
		\item $\widetilde{P} = rP - Q$: the terms $rP$
		      (parity $d\bmod 2$) and $Q$ (parity
		      $(d-2)\bmod 2 = d\bmod 2$) have the
		      \emph{same} parity, so
		      $\widetilde{P}$ has parity $d\bmod 2$.
		\item $\widetilde{Q} = P + rQ$: the terms $P$
		      (parity $(d-1)\bmod 2$) and $rQ$ (parity
		      $(d-1)\bmod 2$) also share the same parity,
		      so $\widetilde{Q}$ has parity
		      $(d-1)\bmod 2$.
	\end{itemize}
	This completes the induction.
\end{proof}

\subsection*{Proof of Eq.~\eqref{eq:error-bound} (Error Amplification)}

\begin{proof}
	The stereographic decoding formula recovers the encoded
	value from Pauli expectation values of the output qubit
	state.  For a real signal, the encoded value is
	\begin{equation}
		z = \frac{\langle\PauliX\rangle}{1 - \langle\PauliZ\rangle},
		\tag{A16}
	\end{equation}
	where $\langle\PauliX\rangle$ and
	$\langle\PauliZ\rangle$ are the Bloch sphere
	coordinates of the state.  For a state encoding the
	real number $r$, the Bloch sphere coordinates are
	\begin{equation}
		\langle\PauliX\rangle = \frac{2r}{1+r^2},
		\qquad
		\langle\PauliZ\rangle = \frac{r^2 - 1}{1+r^2},
		\tag{A17}
	\end{equation}
	which gives
	$1 - \langle\PauliZ\rangle = 1 - (r^2-1)/(1+r^2)
		= 2/(1+r^2)$, confirming that~(A16) returns
	$z = [2r/(1+r^2)]/[2/(1+r^2)] = r$.

	In practice, the expectation values are estimated from
	finite measurements.  Let $\hat{x}$ and $\hat{w}$
	denote the measured estimates of
	$\langle\PauliX\rangle$ and $\langle\PauliZ\rangle$,
	with errors $\delta_x = \hat{x} - \langle\PauliX\rangle$
	and $\delta_w = \hat{w} - \langle\PauliZ\rangle$ each
	bounded by $|\delta_x|, |\delta_w| \leq \delta$.

	The estimated decoded value is
	$\hat{r} = \hat{x}/(1 - \hat{w})$.  To first order in
	$\delta_x, \delta_w$, the error in $\hat{r}$ is
	\begin{equation}
		\hat{r} - r \approx
		\frac{\partial z}{\partial\langle\PauliX\rangle}\,\delta_x
		+ \frac{\partial z}{\partial\langle\PauliZ\rangle}\,\delta_w.
		\tag{A18}
	\end{equation}
	Computing the partial derivatives from~(A16):
	\begin{align}
		\frac{\partial z}{\partial\langle\PauliX\rangle}
		 & = \frac{1}{1-\langle\PauliZ\rangle}
		= \frac{1}{2/(1+r^2)}
		= \frac{1+r^2}{2},
		\tag{A19a}                                                     \\[6pt]
		\frac{\partial z}{\partial\langle\PauliZ\rangle}
		 & = \frac{\langle\PauliX\rangle}{(1-\langle\PauliZ\rangle)^2}
		= \frac{2r/(1+r^2)}{[2/(1+r^2)]^2}
		= \frac{r(1+r^2)}{2}.
		\tag{A19b}
	\end{align}
	The first derivative~(A19a) is computed by
	differentiating $z = x/(1-z_0)$ with respect to~$x$
	at fixed $z_0$; the second~(A19b) uses the quotient
	rule:
	$\partial_{z_0}[x/(1-z_0)] = x/(1-z_0)^2$.

	Applying the triangle inequality to~(A18):
	\begin{align}
		|\hat{r} - r|
		 & \leq \left|\frac{\partial z}{\partial\langle\PauliX\rangle}\right||\delta_x|
		+ \left|\frac{\partial z}{\partial\langle\PauliZ\rangle}\right||\delta_w|
		\nonumber                                                                       \\
		 & \leq \frac{1+r^2}{2}\,\delta + \frac{r(1+r^2)}{2}\,\delta
		\nonumber                                                                       \\
		 & = \frac{(1+r^2)(1+r)}{2}\,\delta.
		\tag{A20}
	\end{align}
	This is~\eqref{eq:error-bound}.  The two contributions
	have distinct physical origins: the first term,
	$(1+r^2)/2$, comes from the sensitivity of the
	numerator measurement and grows as $r^2$ for large~$r$;
	the second term, $r(1+r^2)/2$, comes from the
	denominator $1 - \langle\PauliZ\rangle$ approaching
	zero as $r \to \infty$ (the state approaches the north
	pole of the Bloch sphere) and grows as $r^3$.  The
	denominator contribution dominates for large~$r$,
	reflecting the fundamental geometric fact that the
	stereographic projection compresses the real line near
	infinity into a vanishing arc of the Bloch sphere.
\end{proof}


\subsection*{Proof of Theorem~\ref{thm:gqsp-stereo} (Generalized QSP in Stereo)}

\begin{proof}
	Let $z \in \C\mathbb{P}^1$ and define the unnormalized
	signal unitary
	\begin{equation}
		O_z = \begin{pmatrix} z & -1 \\ 1 & z^* \end{pmatrix},
		\quad \det O_z = 1 + |z|^2,
		\tag{A21}
	\end{equation}
	with $\SU(2)$ control unitaries
	$G_j = \bigl(\begin{smallmatrix} a_j & -b_j^* \\
				b_j & a_j^* \end{smallmatrix}\bigr)$,
	$|a_j|^2 + |b_j|^2 = 1$.
	We prove that there exist $G_0, \ldots, G_d \in \SU(2)$
	such that
	\begin{equation}
		G_d \prod_{j=1}^{d} O_z\,G_j
		= \begin{pmatrix} P(z, z^*) & -Q^*(z, z^*) \\
			Q(z, z^*) & P^*(z, z^*) \end{pmatrix}
		\tag{A22}
	\end{equation}
	if and only if the polynomials $P$ and $Q$ in $z$ and
	$z^*$ satisfy:
	\begin{enumerate}
		\item For all monomials $z^j (z^*)^k$ appearing in
		      $P$ and $Q$, the bi-degree satisfies $j + k \leq d$.
		\item $P(z, z^*)(P(z, z^*))^* + Q(z, z^*)(Q(z, z^*))^*
			      = (1 + |z|^2)^d$.
	\end{enumerate}
	The proof proceeds by induction on the polynomial degree.

	\medskip
	\noindent\textbf{Forward direction.}\quad
	Define the product matrix
	$M_d(z, z^*) = G_d\,O_z\,G_{d-1}\,O_z \cdots G_1\,O_z\,G_0$.
	The base case $M_0 = G_0$ is immediate: $P = a_0$,
	$Q = b_0$, the bi-degree condition holds vacuously,
	and $|a_0|^2 + |b_0|^2 = 1 = (1+|z|^2)^0$.

	For the inductive step, suppose $M_{d-1}$ has the
	claimed form with polynomials $P, Q$ of bi-degree
	$\leq d-1$ satisfying
	$|P|^2 + |Q|^2 = (1+|z|^2)^{d-1}$.  Multiplying by
	$O_z$ from the left:
	\begin{align}
		O_z\,M_{d-1}
		&= \begin{pmatrix} z & -1 \\ 1 & z^* \end{pmatrix}
		\begin{pmatrix} P & -Q^* \\ Q & P^* \end{pmatrix}
		\nonumber\\
		&= \begin{pmatrix} zP - Q & -P^* - zQ^* \\
			P + z^*Q & z^*P^* - Q^* \end{pmatrix}.
		\tag{A23}
	\end{align}
	Applying $G_d$ from the left:
	\begin{align}
		G_d\,O_z\,M_{d-1}
		 & = \begin{pmatrix} a_d & -b_d^* \\
			     b_d & a_d^* \end{pmatrix}
		\begin{pmatrix} zP - Q & -P^* - zQ^* \\
			P + z^*Q & z^*P^* - Q^* \end{pmatrix}
		\nonumber                                                    \\
		 & \equiv \begin{pmatrix} \widetilde{P} & -\widetilde{Q}^* \\
			           \widetilde{Q} & \widetilde{P}^* \end{pmatrix},
		\tag{A24}
	\end{align}
	where
	$\widetilde{P} = a_d(zP - Q) - b_d^*(P + z^*Q)$ and
	$\widetilde{Q} = a_d^*(P + z^*Q) + b_d(zP - Q)$.
	Since $P$ and $Q$ have bi-degree $\leq d-1$,
	the polynomials $\widetilde{P}$ and $\widetilde{Q}$
	have bi-degree $\leq d$.

	To verify the norm condition:
	\begin{align}
		|\widetilde{P}|^2 + |\widetilde{Q}|^2
		 & = (|a_d|^2 + |b_d|^2)\bigl(|zP - Q|^2 + |P + z^*Q|^2\bigr)
		\nonumber                                                        \\
		 & = |zP - Q|^2 + |P + z^*Q|^2.
		\tag{A25}
	\end{align}
	Expanding and using the cancellation of cross terms
	(cf.\ the orthogonality of the columns of~$O_z$):
	\begin{align}
		|zP - Q|^2 + |P + z^*Q|^2
		 & = |z|^2|P|^2 + |Q|^2 + |P|^2 + |z|^2|Q|^2
		\nonumber                                       \\
		 & = (1 + |z|^2)(|P|^2 + |Q|^2)
		\nonumber                                       \\
		 & = (1 + |z|^2)^d.
		\tag{A26}
	\end{align}
	This completes the forward direction.

	\medskip
	\noindent\textbf{Backward direction.}\quad
	Assume the unitary $U$ has the form~(A22) with
	polynomials $P, Q$ of bi-degree $\leq d$ satisfying
	conditions~1 and~2.  We must exhibit $\SU(2)$ controls
	$G_0, \ldots, G_d$ that produce this factorization.
	We proceed by ``peeling off'' one layer at a time.

	Applying $O_z^{-1}\,G_d^{-1}$ to $U$, we obtain
	\begin{equation}
		O_z^{-1}\,G_d^{-1}\,U
		= \begin{pmatrix}
			z^*\widetilde{P} + \widetilde{Q}   &
			-z^*\widetilde{Q}^* + \widetilde{P}^*   \\
			-\widetilde{P} + z\widetilde{Q}     &
			z\widetilde{P}^* + \widetilde{Q}^*
		\end{pmatrix},
		\tag{A27}
	\end{equation}
	where $\widetilde{P}$ and $\widetilde{Q}$ denote the
	entries of $G_d^{-1}\,U$.
	For this to define a valid reduced problem of degree
	$d - 1$, the polynomials
	\begin{equation}
		P'(z,z^*) = \frac{z^*\widetilde{P} + \widetilde{Q}}{1+zz^*},
		\qquad
		Q'(z,z^*) = \frac{z\widetilde{Q} - \widetilde{P}}{1+zz^*}
		\tag{A28}
	\end{equation}
	must be well-defined elements of $\C[z, z^*]$, i.e.,
	$(1+zz^*)$ must divide the numerators.  This
	divisibility is equivalent to requiring that the
	numerators vanish at $z^* = -1/z$, i.e.,
	$P'(z, -1/z) = Q'(z, -1/z) = 0$.

	Enforcing these conditions determines the parameters
	$(a_d, b_d)$ of $G_d$.  Define the Laurent polynomials
	$\hat{p}(z) = P(z, -1/z)$ and
	$\hat{q}(z) = Q(z, -1/z)$.  The normalization
	condition~2 becomes
	\begin{equation}
		\hat{p}(z)\,\hat{p}^*(z) = -\hat{q}(z)\,\hat{q}^*(z).
		\tag{A29}
	\end{equation}
	Since $\C[z, z^{-1}]$ is a unique factorization domain,
	$\hat{p}$ and $\hat{q}$ have a greatest common divisor.
	Let $A = \gcd(\hat{p}, \hat{q})$ and write
	$\hat{p} = AB$, $\hat{q} = AC$ with
	$\gcd(B, C) = 1$.  Then~(A29) becomes
	\begin{equation}
		BB^* = -CC^*.
		\tag{A30}
	\end{equation}
	Since $\gcd(B, C) = 1$, we must have $B \mid C^*$ and
	$C \mid B^*$, which forces $\deg(B) = \deg(C)$.

	\textit{Lemma.}
		If $D \in \C[z, z^{-1}]$ satisfies $DD^* = -1$,
		then $D = cz^k$ with $|c| = 1$ and $k$ odd.

	\begin{proof}[Proof of Lemma]
		Write $D = \sum_n d_n z^n$.  Since $DD^* = -1$,
		the coefficient of $z^0$ in $DD^*$ must satisfy
		$\sum_n |d_n|^2(-1)^n = -1$.  For any power
		$m \neq 0$, the coefficient of $z^m$ is
		$\sum_n d_n d_{n-m}^* (-1)^{n-m} = 0$.
		If two coefficients $d_j, d_k$ with $j \neq k$ are
		nonzero, the coefficient for the $z^{j-k}$ term gives
		$d_j d_k^*(-1)^k = 0$, a contradiction.
		So $D$ has exactly one nonzero term: $D = cz^k$.
		Since $DD^* = |c|^2(-1)^k = -1$, we must have $k$
		odd and $|c| = 1$.
	\end{proof}

	Writing $C^* = BD$ for some $D$, the relation
	$BB^* = -CC^*$ gives $DD^* = -1$, so by Lemma~1,
	$D = cz^k$ with $|c| = 1$ and $k$ odd.  This yields
	two cases.

	\textit{Case 1: $B$ and $C$ are units.}
	Then $\hat{p}/\hat{q} = B/C = cz^k$ for some
	$|c| = 1$ and $k$ odd, and
	$(B/C)(B/C)^* = -1$.  The constraint that $k$ is odd
	forces $k = \pm 1$.  We seek $a, b$ with
	$|a|^2 + |b|^2 = 1$ such that
	$R(z) := (bz + a^*)\hat{p}(z) - (az - b^*)\hat{q}(z)
		\equiv 0$,
	ensuring that $(1 + zz^*)$ divides the numerators
	in~(A28).
	For $k = 1$: the condition $R(z) \equiv 0$ reduces to
	$b = 0$ and $a/a^* = c$, which is satisfiable.
	For $k = -1$: similarly $a = 0$ and
	$-b^*/b = c$, also satisfiable since $|c| = 1$.

	\textit{Case 2: $B$ and $C$ have degree 1.}
	Write $C = (a^* + bz)/\gamma$ and
	$B = (az - b^*)/\delta$ for some $\gamma, \delta$.
	The divisibility condition
	$(bz + a^*)\hat{p}(z) = (az - b^*)\hat{q}(z)$
	requires $\gamma = \delta$, and $B/C = (az-b^*)/(bz+a^*)$.
	Then $R = (bz+a^*)B - (az-b^*)C = 0$.

	In both cases, the $\SU(2)$ parameters are determined:
	$(a, b) \propto (p_d, q_d)$ where $p_d, q_d$ are the
	leading coefficients of $\hat{p}$ and $\hat{q}$,
	normalized so that $|a|^2 + |b|^2 = 1$.  The consistency
	of the $z^{-(d+1)}$ coefficient follows from the
	$z^{2d}$ coefficient of
	$\hat{p}\hat{p}^* + \hat{q}\hat{q}^* = 0$, which gives
	$p_d p_{-d}^* + q_d q_{-d}^* = 0$.

	With $G_d$ so determined, the reduced polynomials
	\begin{align}
		P_r(z, z^*) &= \frac{z^*\widetilde{P} + \widetilde{Q}}{1 + zz^*},
		\nonumber\\
		Q_r(z, z^*) &= \frac{z\widetilde{Q} - \widetilde{P}}{1 + zz^*}
		\tag{A31}
	\end{align}
	are well-defined elements of $\C[z, z^*]$.  It remains
	to verify that the reduced problem satisfies the two
	conditions at degree $d - 1$.  For condition~2:
	\begin{align}
		&|P_r|^2 + |Q_r|^2
		\nonumber\\
		&\quad = \frac{|z^*\widetilde{P} + \widetilde{Q}|^2
		+ |z\widetilde{Q} - \widetilde{P}|^2}{(1+zz^*)^2}
		\nonumber                                          \\
		&\quad = \frac{(1+zz^*)(|\widetilde{P}|^2
		+ |\widetilde{Q}|^2)}{(1+zz^*)^2}
		= (1+|z|^2)^{d-1}.
		\tag{A32}
	\end{align}
	For condition~1, the highest total-degree monomial in
	$(1+|z|^2)^{d-1}$ is $(zz^*)^{d-1}$, which has
	$j + k = 2(d-1)$.  On the left, the highest total
	degree of any monomial $z^j(z^*)^k$ in $P_r$ is
	$2\max\{j+k\}$; for the identity to hold, we need
	$\max(j+k) \leq d-1$.

	Since we can continue peeling in this manner until
	$\deg(P) = 0$ and $\deg(Q) = 0$ (at which point
	$G_0$ is determined directly), this closes the
	induction.
\end{proof}
